\chapter{Память как система: первый обзор}
\label{ch:architecture}

\section{Зачем нужна отдельная карта системы}

Прежде чем обсуждать отдельные механизмы, полезно увидеть всю систему как один
замкнутый цикл. Иначе почти неизбежно возникает ложное ощущение, что память
состоит из набора независимых блоков: где-то отдельно лежит граф, где-то
отдельно работает retrieval, где-то отдельно находится настроение, а языковая
модель ``просто забирает контекст''. Для архитектуры памяти это опасная
иллюзия. Почти каждый блок влияет на остальные двояко: он не только читает
состояние системы, но и перестраивает его.

\section{Что именно строится}

Рассматривается локальный AI-компаньон, память которого должна обладать
следующими свойствами:

\begin{itemize}
\item быть долговременной, а не ограниченной текущим окном контекста;
\item быть ассоциативной, а не только хронологической;
\item различать наблюдение, гипотезу, предпочтение, реактивный паттерн и состояние отношений;
\item допускать пересмотр убеждений;
\item учитывать эмоциональную и личностную окраску того, что всплывает;
\item поддерживать фоновую переработку памяти, а не только обслуживание очередного запроса.
\end{itemize}

Чтобы эти свойства не рассыпались по разным хранилищам и разным схемам
данных, вводится единый субстрат.

\begin{principle}[Единый субстрат]
Память Рин моделируется как один типизированный темпоральный граф с
метаданными на узлах и рёбрах. Все основные процессы памяти либо
непосредственно работают на этом графе, либо сцеплены с ним так, что их
состояние можно выразить через чтение и запись к его полям.
\end{principle}

\section{Четыре слоя системы}

Если смотреть на систему сверху, в ней удобно различать четыре слоя.

\begin{enumerate}
\item Субстрат. Сам граф памяти: узлы, рёбра, уверенность, провенанс, валентность, временной слой, важность и другие поля (глава 5).
\item Операторы. Обход, retrieval, fusion, консолидация, пропагация, забывание, сон и скраб. Это динамика памяти как таковой (главы 6, 10, 11, 14, 12).
\item Состояния. Настроение, эпистемические статусы, вспомогательная метакогнитивная динамика. Эти сущности не следует путать с графом, но они двусторонне с ним связаны (главы 7, 13).
\item Выражение. Слой, который превращает внутреннюю структуру в реплику, тон, фокус внимания и устойчивый характер (главы 16, 17).
\end{enumerate}

\begin{center}
\begin{tikzpicture}[font=\small,>=Stealth,
  layer/.style={draw,rounded corners,minimum width=82mm,align=center,inner sep=6pt,fill=gray!5},
  lbl/.style={font=\bfseries}]
  \node[layer] (expr) {{\bfseries Выражение}\\[1pt]
     \footnotesize личность \quad$\cdot$\quad голос \quad$\cdot$\quad языковая композиция};
  \node[layer,below=7mm of expr] (state) {{\bfseries Состояния}\\[1pt]
     \footnotesize настроение \quad$\cdot$\quad эпистемические статусы \quad$\cdot$\quad метакогнитивные сигналы};
  \node[layer,below=7mm of state] (op) {{\bfseries Операторы}\\[1pt]
     \footnotesize retrieval $\cdot$ обход $\cdot$ fusion $\cdot$ консолидация\\
     \footnotesize пропагация $\cdot$ забывание $\cdot$ сон};
  \node[layer,below=7mm of op] (sub) {{\bfseries Субстрат}\\[1pt]
     \footnotesize граф памяти: узлы, рёбра, типы, метаданные};
  \draw[<->,thick] (expr)--(state);
  \draw[<->,thick] (state)--(op);
  \draw[<->,thick] (op)--(sub);
\end{tikzpicture}
\captionof{figure}{Общая карта системы. Стрелки двусторонние: слой не просто потребляет нижележащий, но и изменяет его состояние.}
\label{fig:layers}
\end{center}

\begin{pitfall}[Чего нельзя делать с этой схемой]
Эти четыре слоя нельзя понимать как четыре независимые базы или четыре
микросервиса с отдельными картинами мира. Слой выражения не должен
придумывать собственную память поверх графа. Слой состояний не должен
дублировать содержательное знание в виде параллельного скрытого графа.
Операторы не должны тайно хранить свою ``локальную истину'', не
согласованную с общим субстратом.
\end{pitfall}

\begin{pitfall}[``Один субстрат'' --- дисциплина, а не догма ``всё есть граф'']
Тезис о едином субстрате легко перетянуть слишком далеко и начать складывать
в граф вообще всё: runtime-очереди, streaming-флаги, транзиентные тики
настроения, состояние планировщика. Это превращает граф не в память, а в
свалку состояний.

\[
\begin{gathered}
\text{граф} = \text{память}, \qquad
\text{event log} = \text{история},\\
\text{телеметрия} = \text{самонаблюдение}, \qquad
\text{scheduler state} = \text{управление},\\
\text{prompt context} = \text{временная сборка для LLM}.
\end{gathered}
\]

Эти сущности связаны, но это разные хранилища с разными законами жизни.
Настроение --- сцепленный сбоку модуль; самотелеметрия метакогниции в
инженерном варианте естественно оформляется как кольцевой буфер; планировщик
--- как конечный автомат режимов работы. Правило простое: если данное
транзиентно или относится к управлению, истории или упаковке контекста, то
оно не должно механически жить в графе памяти.
\end{pitfall}

\section{Что течёт по системе}

Через эту архитектуру проходят несколько разных ``величин'', и их полезно не
смешивать.

\begin{definition}[Активация]
Активация --- это краткоживущая подсветка узла: мера того, насколько он важен
сейчас, в данном локальном ходе памяти.
\end{definition}

\begin{definition}[Уверенность]
Уверенность --- это медленно меняющаяся степень убеждённости в содержании
узла. Она не должна путаться с активацией: можно быть уверенным в факте и при
этом не держать его в фокусе; можно, наоборот, временно активировать
сомнительную гипотезу.
\end{definition}

\begin{definition}[Валентность]
Валентность --- аффективная окраска содержания. Она отвечает не за истинность
узла, а за то, как этот узел переживается и какую роль играет в настроенческой
динамике.
\end{definition}

\begin{definition}[Структурное изменение]
Структурное изменение --- это рождение, укрепление, ослабление, слияние или
удаление узлов и рёбер. В отличие от активации, оно меняет саму геометрию
памяти.
\end{definition}

\section{Сквозной пример}

Полезно проследить полный цикл на одном эпизоде. Пусть во время стрима
зритель пишет: ``Опять слил ранг в доте''.

\begin{workedbox}[От реплики до перестройки памяти]
\begin{enumerate}
\item Кодирование события. В систему поступает наблюдение о конкретном эпизоде. Оно получает тип, базовую уверенность, временной слой и привязку к участнику взаимодействия.
\item Поиск якорей. По содержанию и контексту сцены находятся ближайшие узлы: прежние упоминания доты, прошлые реплики этого зрителя, возможные эмоциональные реакции на проигрыш.
\item Расширение окрестности. От якорей запускается ограниченный обход. Он поднимает не только прямые совпадения, но и соседние ассоциативные структуры.
\item Локальный вывод. Совместная активация нескольких узлов может породить новую гипотезу: например, что в данной сцене уместен мягкий подкол, а не сухое сочувствие.
\item Выражение. На сцену выходит слой личности и голоса: из всплывшего материала строится конкретная реплика в характерном стиле.
\item Обратная запись. Реакция зрителя становится частью новых данных. Если подкол сработал, укрепляется соответствующий паттерн; если нет, появляется основание для пересмотра.
\item Фоновая переработка. Позднее эта сцена может войти в более устойчивое знание: ``данный зритель часто шутит о рейтинге'', ``такая интонация обычно безопасна'', ``эти две темы часто связаны''.
\end{enumerate}
\end{workedbox}

Этот пример показывает главный архитектурный факт: реплика является не
конечной целью отдельного генератора, а одним моментом внутри цикла памяти.
Система не просто говорит; она одновременно перестраивает себя.

\section{Что не входит в граф}

Тезис о едином субстрате не означает, что ``всё вообще'' надо хранить в
памяти-графе.

\begin{itemize}
\item Граф памяти хранит долговременное содержательное знание и его обоснования.
\item Event log хранит историю коммитов и удалений.
\item Телеметрия хранит самонаблюдение и служебные сигналы.
\item Планировщик хранит режимы работы и бюджет вычислений.
\item Промпт-контекст хранит временную сборку того, что подаётся в LLM на конкретном шаге.
\end{itemize}

\begin{intuition}[Почему это различение принципиально]
Если event log, граф, телеметрия и prompt context свалить в одну кучу,
получится не богатая память, а неразличимая смесь истории, управления и
содержания. Такой объект неудобен не только технически. Он делает почти
невозможным ответ на вопрос, что именно является памятью, а что лишь следом
работы памяти.
\end{intuition}

\section{Роль архитектурной главы в остальной книге}

Эта глава не заменяет дальнейшие определения и механизмы. Её задача скромнее:
задать правильную оптику. После неё должно быть ясно не ``как реализуется
каждая функция'', а:

\begin{itemize}
\item что в системе считается субстратом;
\item какие процессы над ним текут;
\item почему они не независимы;
\item где проходит граница между памятью, состоянием и выражением.
\end{itemize}

Следующий шаг поэтому естественен: после общей карты нужен словарь основных
объектов. Этой задаче посвящена глава 4.
